\chapter{Особенности реализации алгоритма}

В качестве языка программирования для реализации алгоритмов был выбран язык программирования Python, разработка велась в среде Jupyter Notebook, позволяющей выполнять программу поэтапно, выводя на экран промежуточные результаты вычислений.

Для построения рекомендательной системы использовался датасет игр с платформы Steam, содержащий 200 000 транзакций с играми, купленными пользователями. Так как датасет не содержит данных о пользовательских сессиях, то рекомендации формировались на основе времени, которое пользователи проводили в играх. Оценкой игры считалось время, проведённое игроком в игре, делённое на среднее время, которое игроки проводили в этой игре. Также был задан верхний порог оценки (5), чтобы уравнять между собой высокие положительные оценки.

Реализованная гибридная рекомендательная система использует матричную факторизацию (коллаборативную фильтрацию) и фильтрацию на основе контента (нахождение похожих игр путём их сравнения по косинусной мере). Оценки, построенные каждым рекомендательным алгоритмом, объединяются с весом $\alpha$ (доля оценок, полученных с помощью матричной факторизации).

\clearpage
